%%%%%%%%%%%%%%%%%%%%%%%%%%%%%%%%%%%%%%%%%%%%%%%%%%%%%%%%%%%%
% 8.16 SPATIAL STATISTICS
% The Composition of Advanced Mathematics
%%%%%%%%%%%%%%%%%%%%%%%%%%%%%%%%%%%%%%%%%%%%%%%%%%%%%%%%%%%%

\section{Spatial Statistics}
\label{sec:spatial-statistics}

\prerequisite{
Probability, random variables, covariance and correlation, regression,
multivariate statistics, hypothesis testing, stochastic processes,
matrix algebra, and basic geometry.
}

\learningobjectives{
By the end of this section, the reader should be able to:
\begin{itemize}
    \item explain what distinguishes spatial data from ordinary
          independent observations;
    \item distinguish geostatistical, areal, and point-pattern data;
    \item define spatial location, distance, and neighborhood structure;
    \item explain spatial dependence and spatial autocorrelation;
    \item understand first-order and second-order spatial structure;
    \item define spatial covariance functions;
    \item understand isotropy and anisotropy;
    \item understand stationarity in spatial processes;
    \item define semivariograms and variograms;
    \item interpret nugget, sill, and range;
    \item construct empirical semivariograms conceptually;
    \item understand common variogram models;
    \item understand kriging conceptually;
    \item distinguish ordinary, simple, and universal kriging;
    \item derive basic kriging weights conceptually;
    \item understand spatial prediction uncertainty;
    \item define spatial weights matrices;
    \item compute and interpret Moran's \(I\);
    \item understand Geary's \(C\);
    \item understand local indicators of spatial association;
    \item distinguish global and local spatial statistics;
    \item understand spatial regression conceptually;
    \item distinguish spatial lag and spatial error models;
    \item understand neighborhood-based areal models;
    \item understand Gaussian random fields conceptually;
    \item understand conditional autoregressive models conceptually;
    \item understand simultaneous autoregressive models conceptually;
    \item understand point processes;
    \item define intensity functions;
    \item understand homogeneous Poisson point processes;
    \item understand Ripley's \(K\)-function conceptually;
    \item distinguish clustering, inhibition, and complete spatial
          randomness;
    \item understand edge effects;
    \item recognize spatial confounding;
    \item understand spatial sampling design;
    \item understand spatial interpolation;
    \item recognize spatiotemporal extensions;
    \item connect spatial statistics with stochastic processes,
          regression, environmental modeling, and statistical learning.
\end{itemize}
}

%%%%%%%%%%%%%%%%%%%%%%%%%%%%%%%%%%%%%%%%%%%%%%%%%%%%%%%%%%%%
\subsection{What Is Spatial Statistics?}
\label{subsec:what-is-spatial-statistics}
%%%%%%%%%%%%%%%%%%%%%%%%%%%%%%%%%%%%%%%%%%%%%%%%%%%%%%%%%%%%

Spatial statistics studies data indexed by geographic or spatial
location.

An observation may be represented as

\[
\boxed{
Y(\mathbf{s}),
}
\]

where

\[
\mathbf{s}
\in
\mathbb{R}^d
\]

is a spatial location.

For two-dimensional geographic data,

\[
\boxed{
\mathbf{s}
=
\begin{pmatrix}
s_1\\
s_2
\end{pmatrix}.
}
\]

%%%%%%%%%%%%%%%%%%%%%%%%%%%%%%%%%%%%%%%%%%%%%%%%%%%%%%%%%%%%
\subsection{Why Location Matters}
\label{subsec:why-location-matters}
%%%%%%%%%%%%%%%%%%%%%%%%%%%%%%%%%%%%%%%%%%%%%%%%%%%%%%%%%%%%

Nearby observations are often more similar than distant observations.

For example, measurements of

\[
\boxed{
\text{temperature},
}
\]

\[
\boxed{
\text{rainfall},
}
\]

\[
\boxed{
\text{soil chemistry},
}
\]

or

\[
\boxed{
\text{disease prevalence}
}
\]

may exhibit spatial dependence.

Thus ordinary independence assumptions may fail.

%%%%%%%%%%%%%%%%%%%%%%%%%%%%%%%%%%%%%%%%%%%%%%%%%%%%%%%%%%%%
\subsection{Spatial Dependence}
\label{subsec:spatial-dependence}
%%%%%%%%%%%%%%%%%%%%%%%%%%%%%%%%%%%%%%%%%%%%%%%%%%%%%%%%%%%%

Spatial dependence means that the distribution of an observation depends
on the values observed at other locations.

Informally,

\[
\boxed{
\text{nearby locations often provide related information}.
}
\]

This dependence may arise from:

\[
\boxed{
\text{physical diffusion},
}
\]

\[
\boxed{
\text{shared environment},
}
\]

\[
\boxed{
\text{movement},
}
\]

\[
\boxed{
\text{social interaction},
}
\]

or

\[
\boxed{
\text{common latent processes}.
}
\]

%%%%%%%%%%%%%%%%%%%%%%%%%%%%%%%%%%%%%%%%%%%%%%%%%%%%%%%%%%%%
\subsection{Three Major Spatial Data Types}
\label{subsec:three-spatial-data-types}
%%%%%%%%%%%%%%%%%%%%%%%%%%%%%%%%%%%%%%%%%%%%%%%%%%%%%%%%%%%%

Spatial statistics commonly distinguishes:

\[
\boxed{
\text{geostatistical data},
}
\]

\[
\boxed{
\text{areal data},
}
\]

and

\[
\boxed{
\text{spatial point patterns}.
}
\]

Each type requires different statistical tools.

%%%%%%%%%%%%%%%%%%%%%%%%%%%%%%%%%%%%%%%%%%%%%%%%%%%%%%%%%%%%
\subsection{Geostatistical Data}
\label{subsec:geostatistical-data}
%%%%%%%%%%%%%%%%%%%%%%%%%%%%%%%%%%%%%%%%%%%%%%%%%%%%%%%%%%%%

Geostatistical data consist of measurements observed at locations within
a continuous spatial region.

Examples include:

\[
\boxed{
\text{temperature at weather stations},
}
\]

\[
\boxed{
\text{dissolved oxygen at monitoring sites},
}
\]

or

\[
\boxed{
\text{soil concentration measurements}.
}
\]

The underlying process is conceptualized as existing throughout the
region, even where it was not observed.

%%%%%%%%%%%%%%%%%%%%%%%%%%%%%%%%%%%%%%%%%%%%%%%%%%%%%%%%%%%%
\subsection{Areal Data}
\label{subsec:areal-data}
%%%%%%%%%%%%%%%%%%%%%%%%%%%%%%%%%%%%%%%%%%%%%%%%%%%%%%%%%%%%

Areal data are aggregated over regions such as:

\[
\boxed{
\text{counties},
}
\]

\[
\boxed{
\text{states},
}
\]

\[
\boxed{
\text{school districts},
}
\]

or

\[
\boxed{
\text{census tracts}.
}
\]

Spatial relationships are often represented through neighborhood or
adjacency structures rather than continuous distance alone.

%%%%%%%%%%%%%%%%%%%%%%%%%%%%%%%%%%%%%%%%%%%%%%%%%%%%%%%%%%%%
\subsection{Spatial Point Patterns}
\label{subsec:spatial-point-patterns}
%%%%%%%%%%%%%%%%%%%%%%%%%%%%%%%%%%%%%%%%%%%%%%%%%%%%%%%%%%%%

Point-pattern data consist of locations where events occur.

Examples include:

\[
\boxed{
\text{tree locations},
}
\]

\[
\boxed{
\text{earthquake epicenters},
}
\]

\[
\boxed{
\text{crime incidents},
}
\]

or

\[
\boxed{
\text{disease cases}.
}
\]

The locations themselves are the primary random objects.

%%%%%%%%%%%%%%%%%%%%%%%%%%%%%%%%%%%%%%%%%%%%%%%%%%%%%%%%%%%%
\subsection{Spatial Coordinates}
\label{subsec:spatial-coordinates}
%%%%%%%%%%%%%%%%%%%%%%%%%%%%%%%%%%%%%%%%%%%%%%%%%%%%%%%%%%%%

A location may be represented by coordinates

\[
\mathbf{s}_i.
\]

For planar data,

\[
\boxed{
\mathbf{s}_i
=
(x_i,y_i).
}
\]

Distance between two locations can be defined using Euclidean distance:

\[
\boxed{
d_{ij}
=
\|
\mathbf{s}_i-\mathbf{s}_j
\|_2.
}
\]

%%%%%%%%%%%%%%%%%%%%%%%%%%%%%%%%%%%%%%%%%%%%%%%%%%%%%%%%%%%%
\subsection{Euclidean Distance}
\label{subsec:spatial-euclidean-distance}
%%%%%%%%%%%%%%%%%%%%%%%%%%%%%%%%%%%%%%%%%%%%%%%%%%%%%%%%%%%%

For

\[
\mathbf{s}_i=(x_i,y_i)
\]

and

\[
\mathbf{s}_j=(x_j,y_j),
\]

the Euclidean distance is

\[
\boxed{
d_{ij}
=
\sqrt{
(x_i-x_j)^2
+
(y_i-y_j)^2
}.
}
\]

This is appropriate when planar geometry is reasonable.

%%%%%%%%%%%%%%%%%%%%%%%%%%%%%%%%%%%%%%%%%%%%%%%%%%%%%%%%%%%%
\subsection{Geographic Distance}
\label{subsec:geographic-distance}
%%%%%%%%%%%%%%%%%%%%%%%%%%%%%%%%%%%%%%%%%%%%%%%%%%%%%%%%%%%%

For large geographic regions, longitude and latitude lie on a curved
surface.

Straight Euclidean distance on raw latitude--longitude coordinates may
therefore be inappropriate.

Possible alternatives include:

\[
\boxed{
\text{great-circle distance},
}
\]

or

\[
\boxed{
\text{projected coordinate systems}.
}
\]

The distance definition should match the spatial scale and geometry.

%%%%%%%%%%%%%%%%%%%%%%%%%%%%%%%%%%%%%%%%%%%%%%%%%%%%%%%%%%%%
\subsection{First-Order Spatial Structure}
\label{subsec:first-order-spatial-structure}
%%%%%%%%%%%%%%%%%%%%%%%%%%%%%%%%%%%%%%%%%%%%%%%%%%%%%%%%%%%%

First-order structure concerns changes in the mean across space.

A spatial process may satisfy

\[
\boxed{
\mathbb{E}
[
Y(\mathbf{s})
]
=
\mu(\mathbf{s}).
}
\]

For example,

\[
\mu(\mathbf{s})
\]

may depend on:

\[
\boxed{
\text{elevation},
}
\]

\[
\boxed{
\text{distance from the coast},
}
\]

or

\[
\boxed{
\text{latitude}.
}
\]

%%%%%%%%%%%%%%%%%%%%%%%%%%%%%%%%%%%%%%%%%%%%%%%%%%%%%%%%%%%%
\subsection{Second-Order Spatial Structure}
\label{subsec:second-order-spatial-structure}
%%%%%%%%%%%%%%%%%%%%%%%%%%%%%%%%%%%%%%%%%%%%%%%%%%%%%%%%%%%%

Second-order structure describes covariance between observations at
different locations.

Define

\[
\boxed{
C(\mathbf{s},\mathbf{t})
=
\operatorname{Cov}
\left[
Y(\mathbf{s}),
Y(\mathbf{t})
\right].
}
\]

This determines how spatial dependence changes with separation.

%%%%%%%%%%%%%%%%%%%%%%%%%%%%%%%%%%%%%%%%%%%%%%%%%%%%%%%%%%%%
\subsection{Second-Order Stationarity}
\label{subsec:spatial-second-order-stationarity}
%%%%%%%%%%%%%%%%%%%%%%%%%%%%%%%%%%%%%%%%%%%%%%%%%%%%%%%%%%%%

A spatial process is second-order stationary if

\[
\boxed{
\mathbb{E}
[
Y(\mathbf{s})
]
=
\mu
}
\]

for every location and

\[
\boxed{
\operatorname{Cov}
[
Y(\mathbf{s}),
Y(\mathbf{s}+\mathbf{h})
]
=
C(\mathbf{h})
}
\]

depends only on the displacement vector

\[
\mathbf{h}.
\]

%%%%%%%%%%%%%%%%%%%%%%%%%%%%%%%%%%%%%%%%%%%%%%%%%%%%%%%%%%%%
\subsection{Spatial Lag Vector}
\label{subsec:spatial-lag-vector}
%%%%%%%%%%%%%%%%%%%%%%%%%%%%%%%%%%%%%%%%%%%%%%%%%%%%%%%%%%%%

For locations

\[
\mathbf{s}
\]

and

\[
\mathbf{s}+\mathbf{h},
\]

the vector

\[
\boxed{
\mathbf{h}
}
\]

is the spatial lag.

It contains both:

\[
\boxed{
\text{distance}
}
\]

and

\[
\boxed{
\text{direction}.
}
\]

%%%%%%%%%%%%%%%%%%%%%%%%%%%%%%%%%%%%%%%%%%%%%%%%%%%%%%%%%%%%
\subsection{Isotropy}
\label{subsec:spatial-isotropy}
%%%%%%%%%%%%%%%%%%%%%%%%%%%%%%%%%%%%%%%%%%%%%%%%%%%%%%%%%%%%

A stationary covariance function is isotropic if dependence depends only
on distance:

\[
\boxed{
C(\mathbf{h})
=
C
(
\|\mathbf{h}\|
).
}
\]

Thus two pairs of locations separated by the same distance have the same
covariance regardless of direction.

%%%%%%%%%%%%%%%%%%%%%%%%%%%%%%%%%%%%%%%%%%%%%%%%%%%%%%%%%%%%
\subsection{Anisotropy}
\label{subsec:spatial-anisotropy}
%%%%%%%%%%%%%%%%%%%%%%%%%%%%%%%%%%%%%%%%%%%%%%%%%%%%%%%%%%%%

A process is anisotropic when spatial dependence differs by direction.

For example, pollution may spread more strongly along prevailing wind
directions.

Then

\[
\boxed{
C(\mathbf{h})
}
\]

depends on both the length and direction of \(\mathbf{h}\).

%%%%%%%%%%%%%%%%%%%%%%%%%%%%%%%%%%%%%%%%%%%%%%%%%%%%%%%%%%%%
\subsection{Spatial Covariance Function}
\label{subsec:spatial-covariance-function}
%%%%%%%%%%%%%%%%%%%%%%%%%%%%%%%%%%%%%%%%%%%%%%%%%%%%%%%%%%%%

A covariance function

\[
\boxed{
C(\mathbf{h})
}
\]

must generate valid positive semidefinite covariance matrices.

For locations

\[
\mathbf{s}_1,\ldots,\mathbf{s}_n,
\]

the matrix

\[
\boxed{
\Sigma_{ij}
=
C
(
\mathbf{s}_i-\mathbf{s}_j
)
}
\]

must satisfy

\[
\boxed{
\mathbf{a}^T
\Sigma
\mathbf{a}
\geq0
}
\]

for every vector \(\mathbf{a}\).

%%%%%%%%%%%%%%%%%%%%%%%%%%%%%%%%%%%%%%%%%%%%%%%%%%%%%%%%%%%%
\subsection{Spatial Correlation Function}
\label{subsec:spatial-correlation-function}
%%%%%%%%%%%%%%%%%%%%%%%%%%%%%%%%%%%%%%%%%%%%%%%%%%%%%%%%%%%%

If

\[
C(\mathbf{0})>0,
\]

the spatial correlation function is

\[
\boxed{
\rho(\mathbf{h})
=
\frac{
C(\mathbf{h})
}{
C(\mathbf{0})
}.
}
\]

For isotropic processes, this is often written

\[
\boxed{
\rho(h),
\qquad
h=\|\mathbf{h}\|.
}
\]

%%%%%%%%%%%%%%%%%%%%%%%%%%%%%%%%%%%%%%%%%%%%%%%%%%%%%%%%%%%%
\subsection{Exponential Covariance Model}
\label{subsec:exponential-spatial-covariance}
%%%%%%%%%%%%%%%%%%%%%%%%%%%%%%%%%%%%%%%%%%%%%%%%%%%%%%%%%%%%

A common isotropic covariance model is

\[
\boxed{
C(h)
=
\sigma^2
e^{-h/\phi},
}
\]

where

\[
\sigma^2
\]

controls variance and

\[
\phi>0
\]

controls spatial correlation length.

As \(h\) increases,

\[
C(h)
\to0.
\]

%%%%%%%%%%%%%%%%%%%%%%%%%%%%%%%%%%%%%%%%%%%%%%%%%%%%%%%%%%%%
\subsection{Gaussian Covariance Model}
\label{subsec:gaussian-spatial-covariance}
%%%%%%%%%%%%%%%%%%%%%%%%%%%%%%%%%%%%%%%%%%%%%%%%%%%%%%%%%%%%

Another model is

\[
\boxed{
C(h)
=
\sigma^2
\exp
\left[
-
\left(
\frac{h}{\phi}
\right)^2
\right].
}
\]

This produces a smoother spatial process than the exponential model.

%%%%%%%%%%%%%%%%%%%%%%%%%%%%%%%%%%%%%%%%%%%%%%%%%%%%%%%%%%%%
\subsection{Matérn Covariance Preview}
\label{subsec:matern-covariance-preview}
%%%%%%%%%%%%%%%%%%%%%%%%%%%%%%%%%%%%%%%%%%%%%%%%%%%%%%%%%%%%

A flexible spatial covariance family is the Matérn family.

One form is

\[
\boxed{
C(h)
=
\sigma^2
\frac{
2^{1-\nu}
}{
\Gamma(\nu)
}
\left(
\frac{h}{\phi}
\right)^\nu
K_\nu
\left(
\frac{h}{\phi}
\right),
}
\]

where \(K_\nu\) is a modified Bessel function.

The parameter

\[
\nu
\]

controls smoothness.

%%%%%%%%%%%%%%%%%%%%%%%%%%%%%%%%%%%%%%%%%%%%%%%%%%%%%%%%%%%%
\subsection{Intrinsic Stationarity}
\label{subsec:intrinsic-stationarity}
%%%%%%%%%%%%%%%%%%%%%%%%%%%%%%%%%%%%%%%%%%%%%%%%%%%%%%%%%%%%

Some spatial processes may not have a stationary covariance function but
their increments may still have stable behavior.

A process is intrinsically stationary if

\[
\boxed{
\mathbb{E}
[
Y(\mathbf{s}+\mathbf{h})
-
Y(\mathbf{s})
]
}
\]

does not depend on \(\mathbf{s}\), and the variance of the increment
depends only on \(\mathbf{h}\).

This motivates variogram methods.

%%%%%%%%%%%%%%%%%%%%%%%%%%%%%%%%%%%%%%%%%%%%%%%%%%%%%%%%%%%%
\subsection{Semivariogram}
\label{subsec:semivariogram}
%%%%%%%%%%%%%%%%%%%%%%%%%%%%%%%%%%%%%%%%%%%%%%%%%%%%%%%%%%%%

The semivariogram is defined by

\[
\boxed{
\gamma(\mathbf{h})
=
\frac12
\operatorname{Var}
\left[
Y(\mathbf{s}+\mathbf{h})
-
Y(\mathbf{s})
\right].
}
\]

Under isotropy,

\[
\boxed{
\gamma(h)
}
\]

depends only on distance \(h\).

%%%%%%%%%%%%%%%%%%%%%%%%%%%%%%%%%%%%%%%%%%%%%%%%%%%%%%%%%%%%
\subsection{Semivariogram and Covariance}
\label{subsec:semivariogram-covariance}
%%%%%%%%%%%%%%%%%%%%%%%%%%%%%%%%%%%%%%%%%%%%%%%%%%%%%%%%%%%%

For a second-order stationary process,

\[
\operatorname{Var}
[
Y(\mathbf{s}+\mathbf{h})
-
Y(\mathbf{s})
]
=
2C(0)-2C(\mathbf{h}).
\]

Therefore,

\[
\boxed{
\gamma(\mathbf{h})
=
C(0)-C(\mathbf{h}).
}
\]

Thus increasing semivariance corresponds to decreasing covariance.

%%%%%%%%%%%%%%%%%%%%%%%%%%%%%%%%%%%%%%%%%%%%%%%%%%%%%%%%%%%%
\subsection{Variogram}
\label{subsec:variogram}
%%%%%%%%%%%%%%%%%%%%%%%%%%%%%%%%%%%%%%%%%%%%%%%%%%%%%%%%%%%%

Some authors call

\[
\boxed{
2\gamma(\mathbf{h})
}
\]

the variogram and

\[
\gamma(\mathbf{h})
\]

the semivariogram.

Terminology varies.

The convention used should therefore be stated explicitly.

%%%%%%%%%%%%%%%%%%%%%%%%%%%%%%%%%%%%%%%%%%%%%%%%%%%%%%%%%%%%
\subsection{Empirical Semivariogram}
\label{subsec:empirical-semivariogram}
%%%%%%%%%%%%%%%%%%%%%%%%%%%%%%%%%%%%%%%%%%%%%%%%%%%%%%%%%%%%

Suppose \(N(h)\) is the set of location pairs whose separation is near
distance \(h\).

A common empirical estimator is

\[
\boxed{
\widehat{\gamma}(h)
=
\frac{
1
}{
2|N(h)|
}
\sum_{(i,j)\in N(h)}
[
Y(\mathbf{s}_i)
-
Y(\mathbf{s}_j)
]^2.
}
\]

Distances are usually grouped into lag bins.

%%%%%%%%%%%%%%%%%%%%%%%%%%%%%%%%%%%%%%%%%%%%%%%%%%%%%%%%%%%%
\subsection{Worked Example: Empirical Semivariance}
\label{subsec:worked-empirical-semivariance}
%%%%%%%%%%%%%%%%%%%%%%%%%%%%%%%%%%%%%%%%%%%%%%%%%%%%%%%%%%%%

\begin{workedexamplebox}{One Distance Bin}

Suppose three pairs of observations at approximately the same separation
have differences

\[
2,
\quad
-4,
\quad
3.
\]

Then

\[
\widehat{\gamma}(h)
=
\frac1{2(3)}
\left[
2^2+(-4)^2+3^2
\right].
\]

Thus,

\[
\begin{answerbox}
\[
\boxed{
\widehat{\gamma}(h)
=
\frac{29}{6}
\approx
4.833.
}
\]
\end{answerbox}

\end{workedexamplebox}

%%%%%%%%%%%%%%%%%%%%%%%%%%%%%%%%%%%%%%%%%%%%%%%%%%%%%%%%%%%%
\subsection{Nugget}
\label{subsec:variogram-nugget}
%%%%%%%%%%%%%%%%%%%%%%%%%%%%%%%%%%%%%%%%%%%%%%%%%%%%%%%%%%%%

The nugget represents semivariance at extremely small spatial separation.

A nonzero nugget may reflect:

\[
\boxed{
\text{measurement error},
}
\]

or

\[
\boxed{
\text{spatial variation below the sampling resolution}.
}
\]

A pure continuous process without measurement error may have

\[
\gamma(0)=0.
\]

%%%%%%%%%%%%%%%%%%%%%%%%%%%%%%%%%%%%%%%%%%%%%%%%%%%%%%%%%%%%
\subsection{Sill}
\label{subsec:variogram-sill}
%%%%%%%%%%%%%%%%%%%%%%%%%%%%%%%%%%%%%%%%%%%%%%%%%%%%%%%%%%%%

The sill is the level approached by the semivariogram at large distances
when spatial correlation becomes negligible.

For a stationary process,

\[
\boxed{
\text{sill}
\approx
C(0)
}
\]

up to nugget conventions.

%%%%%%%%%%%%%%%%%%%%%%%%%%%%%%%%%%%%%%%%%%%%%%%%%%%%%%%%%%%%
\subsection{Range}
\label{subsec:variogram-range}
%%%%%%%%%%%%%%%%%%%%%%%%%%%%%%%%%%%%%%%%%%%%%%%%%%%%%%%%%%%%

The range is the distance beyond which spatial dependence is negligible
under a chosen model.

Some covariance models reach the sill exactly.

Others, such as exponential models, approach it asymptotically.

Then a practical range is often defined using a chosen correlation
threshold.

%%%%%%%%%%%%%%%%%%%%%%%%%%%%%%%%%%%%%%%%%%%%%%%%%%%%%%%%%%%%
\subsection{Spherical Semivariogram Model}
\label{subsec:spherical-semivariogram}
%%%%%%%%%%%%%%%%%%%%%%%%%%%%%%%%%%%%%%%%%%%%%%%%%%%%%%%%%%%%

A common spherical model is

\[
\boxed{
\gamma(h)
=
\begin{cases}
c_0
+
c
\left[
\dfrac{3h}{2a}
-
\dfrac12
\left(
\dfrac{h}{a}
\right)^3
\right],
&
0<h\leq a,
\\[1em]
c_0+c,
&
h>a,
\end{cases}
}
\]

with

\[
\gamma(0)=0
\]

under one common nugget convention.

Here,

\[
c_0
=
\text{nugget},
\]

\[
c_0+c
=
\text{sill},
\]

and

\[
a
=
\text{range}.
\]

%%%%%%%%%%%%%%%%%%%%%%%%%%%%%%%%%%%%%%%%%%%%%%%%%%%%%%%%%%%%
\subsection{Exponential Semivariogram}
\label{subsec:exponential-semivariogram}
%%%%%%%%%%%%%%%%%%%%%%%%%%%%%%%%%%%%%%%%%%%%%%%%%%%%%%%%%%%%

A common exponential semivariogram is

\[
\boxed{
\gamma(h)
=
c_0
+
c
\left[
1-e^{-h/\phi}
\right]
}
\]

for \(h>0\), under a nugget convention.

It approaches the sill

\[
c_0+c
\]

as

\[
h\to\infty.
\]

%%%%%%%%%%%%%%%%%%%%%%%%%%%%%%%%%%%%%%%%%%%%%%%%%%%%%%%%%%%%
\subsection{Directional Variograms}
\label{subsec:directional-variograms}
%%%%%%%%%%%%%%%%%%%%%%%%%%%%%%%%%%%%%%%%%%%%%%%%%%%%%%%%%%%%

To investigate anisotropy, empirical semivariograms can be computed
separately for different directions.

For example:

\[
\boxed{
\text{north--south},
}
\]

\[
\boxed{
\text{east--west},
}
\]

or other angular sectors.

Different ranges or sills by direction suggest anisotropy.

%%%%%%%%%%%%%%%%%%%%%%%%%%%%%%%%%%%%%%%%%%%%%%%%%%%%%%%%%%%%
\subsection{Kriging}
\label{subsec:kriging}
%%%%%%%%%%%%%%%%%%%%%%%%%%%%%%%%%%%%%%%%%%%%%%%%%%%%%%%%%%%%

Kriging is a spatial prediction method that combines observed values
using weights derived from a spatial covariance or variogram model.

A predictor at unobserved location

\[
\mathbf{s}_0
\]

has the form

\[
\boxed{
\widehat{Y}
(
\mathbf{s}_0
)
=
\sum_{i=1}^{n}
\lambda_i
Y
(
\mathbf{s}_i
).
}
\]

The weights depend on spatial dependence.

%%%%%%%%%%%%%%%%%%%%%%%%%%%%%%%%%%%%%%%%%%%%%%%%%%%%%%%%%%%%
\subsection{Kriging Versus Ordinary Weighted Averaging}
\label{subsec:kriging-versus-weighting}
%%%%%%%%%%%%%%%%%%%%%%%%%%%%%%%%%%%%%%%%%%%%%%%%%%%%%%%%%%%%

Kriging does not assign weights solely according to distance.

It accounts for:

\[
\boxed{
\text{distance to the prediction location},
}
\]

and

\[
\boxed{
\text{dependence among observed locations}.
}
\]

Two observations very close to each other may contain redundant
information and therefore receive less combined weight.

%%%%%%%%%%%%%%%%%%%%%%%%%%%%%%%%%%%%%%%%%%%%%%%%%%%%%%%%%%%%
\subsection{Simple Kriging}
\label{subsec:simple-kriging}
%%%%%%%%%%%%%%%%%%%%%%%%%%%%%%%%%%%%%%%%%%%%%%%%%%%%%%%%%%%%

Simple kriging assumes the process mean

\[
\mu
\]

is known.

The predictor may be written

\[
\boxed{
\widehat{Y}(\mathbf{s}_0)
=
\mu
+
\mathbf{c}_0^T
\Sigma^{-1}
(
\mathbf{Y}
-
\mu\mathbf{1}
),
}
\]

where

\[
\mathbf{c}_0
=
\operatorname{Cov}
(
\mathbf{Y},
Y(\mathbf{s}_0)
).
\]

%%%%%%%%%%%%%%%%%%%%%%%%%%%%%%%%%%%%%%%%%%%%%%%%%%%%%%%%%%%%
\subsection{Simple Kriging Weights}
\label{subsec:simple-kriging-weights}
%%%%%%%%%%%%%%%%%%%%%%%%%%%%%%%%%%%%%%%%%%%%%%%%%%%%%%%%%%%%

For a centered process,

\[
\widehat{Y}
(
\mathbf{s}_0
)
=
\boldsymbol{\lambda}^T
\mathbf{Y}.
\]

The minimum mean squared prediction error weights satisfy

\[
\boxed{
\Sigma
\boldsymbol{\lambda}
=
\mathbf{c}_0.
}
\]

Therefore,

\[
\boxed{
\boldsymbol{\lambda}
=
\Sigma^{-1}
\mathbf{c}_0.
}
\]

%%%%%%%%%%%%%%%%%%%%%%%%%%%%%%%%%%%%%%%%%%%%%%%%%%%%%%%%%%%%
\subsection{Simple Kriging Variance}
\label{subsec:simple-kriging-variance}
%%%%%%%%%%%%%%%%%%%%%%%%%%%%%%%%%%%%%%%%%%%%%%%%%%%%%%%%%%%%

Let

\[
C(0)
=
\operatorname{Var}
[
Y(\mathbf{s}_0)
].
\]

The minimum prediction-error variance is

\[
\boxed{
\sigma_K^2
=
C(0)
-
\mathbf{c}_0^T
\Sigma^{-1}
\mathbf{c}_0.
}
\]

This quantifies spatial prediction uncertainty under the assumed model.

%%%%%%%%%%%%%%%%%%%%%%%%%%%%%%%%%%%%%%%%%%%%%%%%%%%%%%%%%%%%
\subsection{Ordinary Kriging}
\label{subsec:ordinary-kriging}
%%%%%%%%%%%%%%%%%%%%%%%%%%%%%%%%%%%%%%%%%%%%%%%%%%%%%%%%%%%%

Ordinary kriging assumes an unknown but constant mean.

The predictor is

\[
\boxed{
\widehat{Y}
(
\mathbf{s}_0
)
=
\sum_{i=1}^{n}
\lambda_i
Y
(
\mathbf{s}_i
),
}
\]

with unbiasedness constraint

\[
\boxed{
\sum_{i=1}^{n}\lambda_i=1.
}
\]

A Lagrange multiplier is used to minimize prediction variance subject to
this constraint.

%%%%%%%%%%%%%%%%%%%%%%%%%%%%%%%%%%%%%%%%%%%%%%%%%%%%%%%%%%%%
\subsection{Ordinary Kriging System}
\label{subsec:ordinary-kriging-system}
%%%%%%%%%%%%%%%%%%%%%%%%%%%%%%%%%%%%%%%%%%%%%%%%%%%%%%%%%%%%

Using covariance notation, the weights satisfy a system of the form

\[
\boxed{
\begin{pmatrix}
\Sigma
&
\mathbf{1}
\\
\mathbf{1}^T
&
0
\end{pmatrix}
\begin{pmatrix}
\boldsymbol{\lambda}\\
\nu
\end{pmatrix}
=
\begin{pmatrix}
\mathbf{c}_0\\
1
\end{pmatrix}.
}
\]

The multiplier

\[
\nu
\]

enforces the unbiasedness constraint.

%%%%%%%%%%%%%%%%%%%%%%%%%%%%%%%%%%%%%%%%%%%%%%%%%%%%%%%%%%%%
\subsection{Universal Kriging}
\label{subsec:universal-kriging}
%%%%%%%%%%%%%%%%%%%%%%%%%%%%%%%%%%%%%%%%%%%%%%%%%%%%%%%%%%%%

Universal kriging allows the mean to vary with spatial covariates.

For example,

\[
\boxed{
\mu(\mathbf{s})
=
\beta_0
+
\beta_1x_1(\mathbf{s})
+
\cdots
+
\beta_px_p(\mathbf{s}).
}
\]

The residual spatial process is modeled through a covariance structure.

Thus universal kriging combines:

\[
\boxed{
\text{regression}
+
\text{spatial covariance}.
}
\]

%%%%%%%%%%%%%%%%%%%%%%%%%%%%%%%%%%%%%%%%%%%%%%%%%%%%%%%%%%%%
\subsection{Regression Kriging}
\label{subsec:regression-kriging}
%%%%%%%%%%%%%%%%%%%%%%%%%%%%%%%%%%%%%%%%%%%%%%%%%%%%%%%%%%%%

A common conceptual decomposition is

\[
\boxed{
Y(\mathbf{s})
=
\mathbf{x}(\mathbf{s})^T
\boldsymbol{\beta}
+
Z(\mathbf{s}),
}
\]

where

\[
Z(\mathbf{s})
\]

is a mean-zero spatially correlated residual process.

Prediction combines the fitted regression component and a kriged
residual.

%%%%%%%%%%%%%%%%%%%%%%%%%%%%%%%%%%%%%%%%%%%%%%%%%%%%%%%%%%%%
\subsection{Gaussian Random Fields}
\label{subsec:gaussian-random-fields}
%%%%%%%%%%%%%%%%%%%%%%%%%%%%%%%%%%%%%%%%%%%%%%%%%%%%%%%%%%%%

A Gaussian random field is a spatial process for which every finite
collection

\[
\boxed{
(
Y(\mathbf{s}_1),
\ldots,
Y(\mathbf{s}_n)
)
}
\]

has a multivariate normal distribution.

Thus the process is completely specified by:

\[
\boxed{
\text{mean function}
}
\]

and

\[
\boxed{
\text{covariance function}.
}
\]

%%%%%%%%%%%%%%%%%%%%%%%%%%%%%%%%%%%%%%%%%%%%%%%%%%%%%%%%%%%%
\subsection{Gaussian Process Spatial Model}
\label{subsec:gaussian-process-spatial}
%%%%%%%%%%%%%%%%%%%%%%%%%%%%%%%%%%%%%%%%%%%%%%%%%%%%%%%%%%%%

A common model is

\[
\boxed{
Y(\mathbf{s})
=
\mu(\mathbf{s})
+
Z(\mathbf{s})
+
\varepsilon(\mathbf{s}),
}
\]

where

\[
Z(\mathbf{s})
\]

is a Gaussian spatial process and

\[
\varepsilon(\mathbf{s})
\]

is independent measurement error.

This naturally produces a nugget effect.

%%%%%%%%%%%%%%%%%%%%%%%%%%%%%%%%%%%%%%%%%%%%%%%%%%%%%%%%%%%%
\subsection{Spatial Prediction as Conditional Normal Theory}
\label{subsec:spatial-conditional-normal}
%%%%%%%%%%%%%%%%%%%%%%%%%%%%%%%%%%%%%%%%%%%%%%%%%%%%%%%%%%%%

Under a Gaussian process,

\[
\begin{pmatrix}
\mathbf{Y}\\
Y(\mathbf{s}_0)
\end{pmatrix}
\]

is jointly normal.

Therefore the conditional distribution

\[
\boxed{
Y(\mathbf{s}_0)
\mid
\mathbf{Y}
}
\]

is normal.

Its conditional mean produces the kriging predictor, while its
conditional variance gives prediction uncertainty.

%%%%%%%%%%%%%%%%%%%%%%%%%%%%%%%%%%%%%%%%%%%%%%%%%%%%%%%%%%%%
\subsection{Spatial Weights Matrix}
\label{subsec:spatial-weights-matrix}
%%%%%%%%%%%%%%%%%%%%%%%%%%%%%%%%%%%%%%%%%%%%%%%%%%%%%%%%%%%%

For areal data, spatial relationships are often represented by a matrix

\[
\boxed{
W=(w_{ij}).
}
\]

The entry

\[
w_{ij}
\]

describes how strongly region \(j\) is considered a neighbor of region
\(i\).

%%%%%%%%%%%%%%%%%%%%%%%%%%%%%%%%%%%%%%%%%%%%%%%%%%%%%%%%%%%%
\subsection{Binary Adjacency Weights}
\label{subsec:binary-adjacency-weights}
%%%%%%%%%%%%%%%%%%%%%%%%%%%%%%%%%%%%%%%%%%%%%%%%%%%%%%%%%%%%

A simple choice is

\[
\boxed{
w_{ij}
=
\begin{cases}
1,
&
\text{regions }i\text{ and }j\text{ are neighbors},
\\
0,
&
\text{otherwise}.
\end{cases}
}
\]

Typically,

\[
\boxed{
w_{ii}=0.
}
\]

%%%%%%%%%%%%%%%%%%%%%%%%%%%%%%%%%%%%%%%%%%%%%%%%%%%%%%%%%%%%
\subsection{Rook and Queen Adjacency}
\label{subsec:rook-queen-adjacency}
%%%%%%%%%%%%%%%%%%%%%%%%%%%%%%%%%%%%%%%%%%%%%%%%%%%%%%%%%%%%

For polygonal regions:

\[
\boxed{
\text{rook adjacency}
}
\]

requires a shared boundary segment.

\[
\boxed{
\text{queen adjacency}
}
\]

allows either a shared boundary or a shared corner.

The chosen neighborhood definition can affect results.

%%%%%%%%%%%%%%%%%%%%%%%%%%%%%%%%%%%%%%%%%%%%%%%%%%%%%%%%%%%%
\subsection{Distance-Based Weights}
\label{subsec:distance-based-spatial-weights}
%%%%%%%%%%%%%%%%%%%%%%%%%%%%%%%%%%%%%%%%%%%%%%%%%%%%%%%%%%%%

Weights may depend on distance.

For example,

\[
\boxed{
w_{ij}
=
\frac1{
d_{ij}
}
}
\]

for \(i\neq j\), or

\[
\boxed{
w_{ij}
=
\mathbf{1}
\{
d_{ij}\leq d_0
\}.
}
\]

The weighting rule should represent the scientific mechanism of spatial
interaction.

%%%%%%%%%%%%%%%%%%%%%%%%%%%%%%%%%%%%%%%%%%%%%%%%%%%%%%%%%%%%
\subsection{Row Standardization}
\label{subsec:row-standardization-spatial}
%%%%%%%%%%%%%%%%%%%%%%%%%%%%%%%%%%%%%%%%%%%%%%%%%%%%%%%%%%%%

A weights matrix is often row-standardized:

\[
\boxed{
w_{ij}^*
=
\frac{
w_{ij}
}{
\sum_jw_{ij}
}.
}
\]

Then each row sums to \(1\).

This makes

\[
\sum_j
w_{ij}^*Y_j
\]

a weighted average of neighboring values.

%%%%%%%%%%%%%%%%%%%%%%%%%%%%%%%%%%%%%%%%%%%%%%%%%%%%%%%%%%%%
\subsection{Spatial Lag}
\label{subsec:spatial-lag}
%%%%%%%%%%%%%%%%%%%%%%%%%%%%%%%%%%%%%%%%%%%%%%%%%%%%%%%%%%%%

For areal observations stored in vector

\[
\mathbf{y},
\]

the spatial lag is

\[
\boxed{
W\mathbf{y}.
}
\]

The \(i\)-th component is

\[
\boxed{
\sum_j
w_{ij}y_j.
}
\]

It summarizes neighboring outcomes.

%%%%%%%%%%%%%%%%%%%%%%%%%%%%%%%%%%%%%%%%%%%%%%%%%%%%%%%%%%%%
\subsection{Spatial Autocorrelation}
\label{subsec:spatial-autocorrelation}
%%%%%%%%%%%%%%%%%%%%%%%%%%%%%%%%%%%%%%%%%%%%%%%%%%%%%%%%%%%%

Spatial autocorrelation measures similarity between values at nearby
locations.

Positive spatial autocorrelation means nearby observations tend to be
similar.

Negative spatial autocorrelation means neighboring observations tend to
be dissimilar.

No spatial autocorrelation means the spatial arrangement carries little
systematic information about the observed values.

%%%%%%%%%%%%%%%%%%%%%%%%%%%%%%%%%%%%%%%%%%%%%%%%%%%%%%%%%%%%
\subsection{Moran's \(I\)}
\label{subsec:morans-i}
%%%%%%%%%%%%%%%%%%%%%%%%%%%%%%%%%%%%%%%%%%%%%%%%%%%%%%%%%%%%

A widely used global spatial autocorrelation statistic is Moran's \(I\):

\[
\boxed{
I
=
\frac{
n
}{
S_0
}
\frac{
\displaystyle
\sum_i
\sum_j
w_{ij}
(
y_i-\overline{y}
)
(
y_j-\overline{y}
)
}{
\displaystyle
\sum_i
(
y_i-\overline{y}
)^2
},
}
\]

where

\[
\boxed{
S_0
=
\sum_i
\sum_j
w_{ij}.
}
\]

%%%%%%%%%%%%%%%%%%%%%%%%%%%%%%%%%%%%%%%%%%%%%%%%%%%%%%%%%%%%
\subsection{Interpretation of Moran's \(I\)}
\label{subsec:morans-i-interpretation}
%%%%%%%%%%%%%%%%%%%%%%%%%%%%%%%%%%%%%%%%%%%%%%%%%%%%%%%%%%%%

Large positive values indicate spatial clustering of similar values.

Negative values indicate neighboring dissimilarity.

The exact expected value under a randomization null is often

\[
\boxed{
\mathbb{E}[I]
=
-\frac1{n-1}
}
\]

under common assumptions.

Therefore the null expectation is not generally exactly zero.

%%%%%%%%%%%%%%%%%%%%%%%%%%%%%%%%%%%%%%%%%%%%%%%%%%%%%%%%%%%%
\subsection{Worked Example: Moran Structure}
\label{subsec:worked-moran-structure}
%%%%%%%%%%%%%%%%%%%%%%%%%%%%%%%%%%%%%%%%%%%%%%%%%%%%%%%%%%%%

\begin{workedexamplebox}{Neighbor Similarity}

Suppose neighboring regions tend to have outcomes both above the mean or
both below the mean.

Then terms

\[
(
y_i-\overline{y}
)
(
y_j-\overline{y}
)
\]

are often positive.

Therefore the numerator of Moran's \(I\) tends to increase.

\begin{answerbox}
\[
\boxed{
I>0
\text{ suggests positive spatial autocorrelation.}
}
\]
\end{answerbox}

\end{workedexamplebox}

%%%%%%%%%%%%%%%%%%%%%%%%%%%%%%%%%%%%%%%%%%%%%%%%%%%%%%%%%%%%
\subsection{Geary's \(C\)}
\label{subsec:gearys-c}
%%%%%%%%%%%%%%%%%%%%%%%%%%%%%%%%%%%%%%%%%%%%%%%%%%%%%%%%%%%%

Another global spatial statistic is Geary's \(C\):

\[
\boxed{
C
=
\frac{
(n-1)
\displaystyle
\sum_i
\sum_j
w_{ij}
(
y_i-y_j
)^2
}{
2S_0
\displaystyle
\sum_i
(
y_i-\overline{y}
)^2
}.
}
\]

Unlike Moran's \(I\), it is based directly on neighboring differences.

%%%%%%%%%%%%%%%%%%%%%%%%%%%%%%%%%%%%%%%%%%%%%%%%%%%%%%%%%%%%
\subsection{Interpretation of Geary's \(C\)}
\label{subsec:gearys-c-interpretation}
%%%%%%%%%%%%%%%%%%%%%%%%%%%%%%%%%%%%%%%%%%%%%%%%%%%%%%%%%%%%

Under common null formulations,

\[
\boxed{
C\approx1
}
\]

indicates little spatial autocorrelation.

Values

\[
\boxed{
C<1
}
\]

suggest positive spatial autocorrelation.

Values

\[
\boxed{
C>1
}
\]

suggest negative spatial autocorrelation.

%%%%%%%%%%%%%%%%%%%%%%%%%%%%%%%%%%%%%%%%%%%%%%%%%%%%%%%%%%%%
\subsection{Global Versus Local Spatial Statistics}
\label{subsec:global-versus-local-spatial}
%%%%%%%%%%%%%%%%%%%%%%%%%%%%%%%%%%%%%%%%%%%%%%%%%%%%%%%%%%%%

Global statistics summarize spatial dependence across the entire study
region.

Local statistics identify spatial patterns around particular locations.

Thus:

\[
\boxed{
\text{global}
=
\text{overall spatial structure},
}
\]

while

\[
\boxed{
\text{local}
=
\text{location-specific spatial structure}.
}
\]

%%%%%%%%%%%%%%%%%%%%%%%%%%%%%%%%%%%%%%%%%%%%%%%%%%%%%%%%%%%%
\subsection{Local Moran's \(I\)}
\label{subsec:local-morans-i}
%%%%%%%%%%%%%%%%%%%%%%%%%%%%%%%%%%%%%%%%%%%%%%%%%%%%%%%%%%%%

A common local Moran statistic has the form

\[
\boxed{
I_i
=
\frac{
y_i-\overline{y}
}{
m_2
}
\sum_j
w_{ij}
(
y_j-\overline{y}
),
}
\]

where \(m_2\) is a variance-like normalization.

It measures whether observation \(i\) is surrounded by similar or
dissimilar values.

%%%%%%%%%%%%%%%%%%%%%%%%%%%%%%%%%%%%%%%%%%%%%%%%%%%%%%%%%%%%
\subsection{LISA}
\label{subsec:lisa}
%%%%%%%%%%%%%%%%%%%%%%%%%%%%%%%%%%%%%%%%%%%%%%%%%%%%%%%%%%%%

Local indicators of spatial association are often abbreviated LISA.

They can identify:

\[
\boxed{
\text{high--high clusters},
}
\]

\[
\boxed{
\text{low--low clusters},
}
\]

\[
\boxed{
\text{high--low outliers},
}
\]

and

\[
\boxed{
\text{low--high outliers}.
}
\]

Multiple-testing issues arise when many local statistics are examined.

%%%%%%%%%%%%%%%%%%%%%%%%%%%%%%%%%%%%%%%%%%%%%%%%%%%%%%%%%%%%
\subsection{Spatial Regression}
\label{subsec:spatial-regression}
%%%%%%%%%%%%%%%%%%%%%%%%%%%%%%%%%%%%%%%%%%%%%%%%%%%%%%%%%%%%

Ordinary regression assumes errors are independent after conditioning on
predictors.

Spatial data may violate this assumption.

A basic regression model is

\[
\boxed{
\mathbf{y}
=
X\boldsymbol{\beta}
+
\boldsymbol{\varepsilon}.
}
\]

If

\[
\boldsymbol{\varepsilon}
\]

is spatially correlated, ordinary least-squares standard errors may be
incorrect.

%%%%%%%%%%%%%%%%%%%%%%%%%%%%%%%%%%%%%%%%%%%%%%%%%%%%%%%%%%%%
\subsection{Spatial Error Model}
\label{subsec:spatial-error-model}
%%%%%%%%%%%%%%%%%%%%%%%%%%%%%%%%%%%%%%%%%%%%%%%%%%%%%%%%%%%%

A spatial error model places spatial dependence in the residual process:

\[
\boxed{
\mathbf{y}
=
X\boldsymbol{\beta}
+
\mathbf{u},
}
\]

with

\[
\boxed{
\mathbf{u}
=
\lambda
W\mathbf{u}
+
\boldsymbol{\varepsilon}.
}
\]

Therefore,

\[
\boxed{
(
I-\lambda W
)
\mathbf{u}
=
\boldsymbol{\varepsilon}.
}
\]

%%%%%%%%%%%%%%%%%%%%%%%%%%%%%%%%%%%%%%%%%%%%%%%%%%%%%%%%%%%%
\subsection{Spatial Lag Model}
\label{subsec:spatial-lag-model}
%%%%%%%%%%%%%%%%%%%%%%%%%%%%%%%%%%%%%%%%%%%%%%%%%%%%%%%%%%%%

A spatial lag model includes neighboring outcomes directly:

\[
\boxed{
\mathbf{y}
=
\rho
W\mathbf{y}
+
X\boldsymbol{\beta}
+
\boldsymbol{\varepsilon}.
}
\]

The parameter

\[
\rho
\]

controls dependence on neighboring response values.

This is also called a spatial autoregressive model in some contexts.

%%%%%%%%%%%%%%%%%%%%%%%%%%%%%%%%%%%%%%%%%%%%%%%%%%%%%%%%%%%%
\subsection{Reduced Form of the Spatial Lag Model}
\label{subsec:spatial-lag-reduced-form}
%%%%%%%%%%%%%%%%%%%%%%%%%%%%%%%%%%%%%%%%%%%%%%%%%%%%%%%%%%%%

Rearrange:

\[
(
I-\rho W
)
\mathbf{y}
=
X\boldsymbol{\beta}
+
\boldsymbol{\varepsilon}.
\]

If the inverse exists,

\[
\boxed{
\mathbf{y}
=
(
I-\rho W
)^{-1}
X\boldsymbol{\beta}
+
(
I-\rho W
)^{-1}
\boldsymbol{\varepsilon}.
}
\]

Thus a local change can propagate through the spatial network.

%%%%%%%%%%%%%%%%%%%%%%%%%%%%%%%%%%%%%%%%%%%%%%%%%%%%%%%%%%%%
\subsection{Spatial Lag Versus Spatial Error Dependence}
\label{subsec:spatial-lag-versus-error}
%%%%%%%%%%%%%%%%%%%%%%%%%%%%%%%%%%%%%%%%%%%%%%%%%%%%%%%%%%%%

A spatial lag model says neighboring outcomes directly enter the response
structure.

A spatial error model says dependence remains in unobserved residual
factors.

These correspond to different scientific mechanisms and should not be
selected solely by whichever model fits better numerically.

%%%%%%%%%%%%%%%%%%%%%%%%%%%%%%%%%%%%%%%%%%%%%%%%%%%%%%%%%%%%
\subsection{Spatial Confounding}
\label{subsec:spatial-confounding}
%%%%%%%%%%%%%%%%%%%%%%%%%%%%%%%%%%%%%%%%%%%%%%%%%%%%%%%%%%%%

Spatial confounding can arise when a predictor has spatial structure that
overlaps strongly with an unobserved spatial effect.

Then estimated covariate effects and spatial random effects may be
difficult to distinguish.

This can increase uncertainty or alter coefficient interpretation.

%%%%%%%%%%%%%%%%%%%%%%%%%%%%%%%%%%%%%%%%%%%%%%%%%%%%%%%%%%%%
\subsection{Conditional Autoregressive Models}
\label{subsec:conditional-autoregressive-models}
%%%%%%%%%%%%%%%%%%%%%%%%%%%%%%%%%%%%%%%%%%%%%%%%%%%%%%%%%%%%

Conditional autoregressive, or CAR, models are commonly used for areal
spatial random effects.

The distribution of spatial effect \(U_i\) is specified conditionally on
neighboring effects.

Conceptually,

\[
\boxed{
U_i
\mid
U_{-i}
}
\]

has a mean related to nearby

\[
U_j.
\]

%%%%%%%%%%%%%%%%%%%%%%%%%%%%%%%%%%%%%%%%%%%%%%%%%%%%%%%%%%%%
\subsection{Simple CAR Structure}
\label{subsec:simple-car-structure}
%%%%%%%%%%%%%%%%%%%%%%%%%%%%%%%%%%%%%%%%%%%%%%%%%%%%%%%%%%%%

A simplified conditional mean may have the form

\[
\boxed{
\mathbb{E}
[
U_i
\mid
U_{-i}
]
=
\frac{
\sum_j
w_{ij}U_j
}{
\sum_jw_{ij}
}.
}
\]

Thus each region tends toward a weighted average of its neighbors.

The corresponding conditional variance often decreases when a region has
more neighbors.

%%%%%%%%%%%%%%%%%%%%%%%%%%%%%%%%%%%%%%%%%%%%%%%%%%%%%%%%%%%%
\subsection{Intrinsic CAR Preview}
\label{subsec:intrinsic-car-preview}
%%%%%%%%%%%%%%%%%%%%%%%%%%%%%%%%%%%%%%%%%%%%%%%%%%%%%%%%%%%%

Intrinsic CAR models use local difference penalties such as

\[
\boxed{
\sum_{i\sim j}
(
U_i-U_j
)^2.
}
\]

These models encourage neighboring effects to be similar.

Some intrinsic CAR priors are improper until an identifiability
constraint such as

\[
\boxed{
\sum_iU_i=0
}
\]

is imposed.

%%%%%%%%%%%%%%%%%%%%%%%%%%%%%%%%%%%%%%%%%%%%%%%%%%%%%%%%%%%%
\subsection{Simultaneous Autoregressive Models}
\label{subsec:simultaneous-autoregressive-models}
%%%%%%%%%%%%%%%%%%%%%%%%%%%%%%%%%%%%%%%%%%%%%%%%%%%%%%%%%%%%

A simultaneous autoregressive, or SAR, model may be written

\[
\boxed{
\mathbf{u}
=
\rho W\mathbf{u}
+
\boldsymbol{\varepsilon}.
}
\]

Therefore,

\[
\boxed{
\mathbf{u}
=
(
I-\rho W
)^{-1}
\boldsymbol{\varepsilon}.
}
\]

This gives an explicit joint covariance structure.

%%%%%%%%%%%%%%%%%%%%%%%%%%%%%%%%%%%%%%%%%%%%%%%%%%%%%%%%%%%%
\subsection{CAR Versus SAR}
\label{subsec:car-versus-sar}
%%%%%%%%%%%%%%%%%%%%%%%%%%%%%%%%%%%%%%%%%%%%%%%%%%%%%%%%%%%%

CAR models are usually specified through conditional distributions.

SAR models are usually specified through simultaneous equations.

Both generate spatial dependence, but their mathematical constructions
differ.

%%%%%%%%%%%%%%%%%%%%%%%%%%%%%%%%%%%%%%%%%%%%%%%%%%%%%%%%%%%%
\subsection{Spatial Generalized Linear Models}
\label{subsec:spatial-generalized-linear-models}
%%%%%%%%%%%%%%%%%%%%%%%%%%%%%%%%%%%%%%%%%%%%%%%%%%%%%%%%%%%%

Spatial dependence can be incorporated into models for non-Gaussian
outcomes.

For binary data,

\[
\boxed{
Y_i
\sim
\operatorname{Bernoulli}(p_i),
}
\]

with

\[
\boxed{
\operatorname{logit}(p_i)
=
\mathbf{x}_i^T
\boldsymbol{\beta}
+
U_i.
}
\]

For count data,

\[
\boxed{
Y_i
\sim
\operatorname{Poisson}(\lambda_i),
}
\]

with

\[
\boxed{
\log\lambda_i
=
\mathbf{x}_i^T
\boldsymbol{\beta}
+
U_i.
}
\]

%%%%%%%%%%%%%%%%%%%%%%%%%%%%%%%%%%%%%%%%%%%%%%%%%%%%%%%%%%%%
\subsection{Point Processes}
\label{subsec:point-processes}
%%%%%%%%%%%%%%%%%%%%%%%%%%%%%%%%%%%%%%%%%%%%%%%%%%%%%%%%%%%%

A spatial point process models a random collection of event locations.

Let

\[
N(A)
\]

denote the number of events in spatial region \(A\).

The process is characterized by probabilities for these random counts and
event configurations.

%%%%%%%%%%%%%%%%%%%%%%%%%%%%%%%%%%%%%%%%%%%%%%%%%%%%%%%%%%%%
\subsection{Intensity Function}
\label{subsec:spatial-intensity-function}
%%%%%%%%%%%%%%%%%%%%%%%%%%%%%%%%%%%%%%%%%%%%%%%%%%%%%%%%%%%%

The intensity function

\[
\boxed{
\lambda(\mathbf{s})
}
\]

describes expected event density near location \(\mathbf{s}\).

Informally,

\[
\boxed{
\lambda(\mathbf{s})
\approx
\frac{
\mathbb{E}
[
N(d\mathbf{s})
]
}{
|d\mathbf{s}|
}.
}
\]

For a homogeneous process,

\[
\boxed{
\lambda(\mathbf{s})
=
\lambda.
}
\]

%%%%%%%%%%%%%%%%%%%%%%%%%%%%%%%%%%%%%%%%%%%%%%%%%%%%%%%%%%%%
\subsection{Homogeneous Poisson Point Process}
\label{subsec:homogeneous-poisson-point-process}
%%%%%%%%%%%%%%%%%%%%%%%%%%%%%%%%%%%%%%%%%%%%%%%%%%%%%%%%%%%%

A homogeneous Poisson point process has constant intensity

\[
\lambda.
\]

For region \(A\),

\[
\boxed{
N(A)
\sim
\operatorname{Poisson}
(
\lambda|A|
),
}
\]

where

\[
|A|
\]

is the area of \(A\).

Counts in disjoint regions are independent.

%%%%%%%%%%%%%%%%%%%%%%%%%%%%%%%%%%%%%%%%%%%%%%%%%%%%%%%%%%%%
\subsection{Complete Spatial Randomness}
\label{subsec:complete-spatial-randomness}
%%%%%%%%%%%%%%%%%%%%%%%%%%%%%%%%%%%%%%%%%%%%%%%%%%%%%%%%%%%%

A homogeneous Poisson point process is commonly used as a model of
complete spatial randomness, abbreviated CSR.

Under CSR:

\[
\boxed{
\text{intensity is constant},
}
\]

and

\[
\boxed{
\text{events do not interact}.
}
\]

Observed departures may suggest clustering or inhibition.

%%%%%%%%%%%%%%%%%%%%%%%%%%%%%%%%%%%%%%%%%%%%%%%%%%%%%%%%%%%%
\subsection{Clustering and Inhibition}
\label{subsec:clustering-inhibition}
%%%%%%%%%%%%%%%%%%%%%%%%%%%%%%%%%%%%%%%%%%%%%%%%%%%%%%%%%%%%

A clustered point pattern contains more nearby event pairs than expected
under CSR.

An inhibited or regular pattern contains fewer nearby pairs.

Examples include:

\[
\boxed{
\text{disease clusters}
}
\]

and

\[
\boxed{
\text{territorial trees or animals}.
}
\]

%%%%%%%%%%%%%%%%%%%%%%%%%%%%%%%%%%%%%%%%%%%%%%%%%%%%%%%%%%%%
\subsection{Nearest-Neighbor Distance}
\label{subsec:nearest-neighbor-distance}
%%%%%%%%%%%%%%%%%%%%%%%%%%%%%%%%%%%%%%%%%%%%%%%%%%%%%%%%%%%%

For event \(i\), define the nearest-neighbor distance

\[
\boxed{
D_i
=
\min_{j\neq i}
\|
\mathbf{s}_i-\mathbf{s}_j
\|.
}
\]

The distribution of these distances provides information about spatial
clustering or inhibition.

%%%%%%%%%%%%%%%%%%%%%%%%%%%%%%%%%%%%%%%%%%%%%%%%%%%%%%%%%%%%
\subsection{Ripley's \(K\)-Function}
\label{subsec:ripleys-k-function}
%%%%%%%%%%%%%%%%%%%%%%%%%%%%%%%%%%%%%%%%%%%%%%%%%%%%%%%%%%%%

Ripley's \(K\)-function summarizes the expected number of other events
within distance \(r\) of a typical event, adjusted for intensity.

Under homogeneous CSR in two dimensions,

\[
\boxed{
K(r)
=
\pi r^2.
}
\]

%%%%%%%%%%%%%%%%%%%%%%%%%%%%%%%%%%%%%%%%%%%%%%%%%%%%%%%%%%%%
\subsection{Interpretation of \(K(r)\)}
\label{subsec:ripleys-k-interpretation}
%%%%%%%%%%%%%%%%%%%%%%%%%%%%%%%%%%%%%%%%%%%%%%%%%%%%%%%%%%%%

If an estimated \(K(r)\) exceeds

\[
\pi r^2,
\]

there may be more nearby event pairs than expected under CSR.

This suggests clustering at scale \(r\).

If it is below

\[
\pi r^2,
\]

the pattern may show inhibition.

Inference should account for sampling variability and edge effects.

%%%%%%%%%%%%%%%%%%%%%%%%%%%%%%%%%%%%%%%%%%%%%%%%%%%%%%%%%%%%
\subsection{\(L\)-Function Transformation}
\label{subsec:ripley-l-function}
%%%%%%%%%%%%%%%%%%%%%%%%%%%%%%%%%%%%%%%%%%%%%%%%%%%%%%%%%%%%

A common transformation is

\[
\boxed{
L(r)
=
\sqrt{
\frac{
K(r)
}{
\pi
}
}.
}
\]

Under CSR,

\[
\boxed{
L(r)=r.
}
\]

Thus plotting

\[
L(r)-r
\]

around zero simplifies interpretation.

%%%%%%%%%%%%%%%%%%%%%%%%%%%%%%%%%%%%%%%%%%%%%%%%%%%%%%%%%%%%
\subsection{Edge Effects}
\label{subsec:spatial-edge-effects}
%%%%%%%%%%%%%%%%%%%%%%%%%%%%%%%%%%%%%%%%%%%%%%%%%%%%%%%%%%%%

Events near the boundary of the observation region have potential
neighbors outside the observed region.

If ignored, this can bias estimates of neighborhood statistics.

This is called an edge effect.

Several correction methods reweight or adjust boundary observations.

%%%%%%%%%%%%%%%%%%%%%%%%%%%%%%%%%%%%%%%%%%%%%%%%%%%%%%%%%%%%
\subsection{Inhomogeneous Point Processes}
\label{subsec:inhomogeneous-point-processes}
%%%%%%%%%%%%%%%%%%%%%%%%%%%%%%%%%%%%%%%%%%%%%%%%%%%%%%%%%%%%

Event intensity may vary across space because of measured environmental
factors.

A model may use

\[
\boxed{
\log
\lambda(\mathbf{s})
=
\beta_0
+
\beta_1x_1(\mathbf{s})
+
\cdots
+
\beta_px_p(\mathbf{s}).
}
\]

Then apparent clustering may be explained partly by changing first-order
intensity.

%%%%%%%%%%%%%%%%%%%%%%%%%%%%%%%%%%%%%%%%%%%%%%%%%%%%%%%%%%%%
\subsection{Cox Processes Preview}
\label{subsec:cox-processes-preview}
%%%%%%%%%%%%%%%%%%%%%%%%%%%%%%%%%%%%%%%%%%%%%%%%%%%%%%%%%%%%

A Cox process is a Poisson point process with a random intensity field.

Conditional on

\[
\lambda(\mathbf{s}),
\]

events follow a Poisson process.

Random variation in the intensity induces extra clustering.

%%%%%%%%%%%%%%%%%%%%%%%%%%%%%%%%%%%%%%%%%%%%%%%%%%%%%%%%%%%%
\subsection{Log-Gaussian Cox Process Preview}
\label{subsec:log-gaussian-cox-process}
%%%%%%%%%%%%%%%%%%%%%%%%%%%%%%%%%%%%%%%%%%%%%%%%%%%%%%%%%%%%

A log-Gaussian Cox process assumes

\[
\boxed{
\log
\lambda(\mathbf{s})
=
\mu(\mathbf{s})
+
Z(\mathbf{s}),
}
\]

where

\[
Z(\mathbf{s})
\]

is a Gaussian random field.

This combines point-process and geostatistical modeling.

%%%%%%%%%%%%%%%%%%%%%%%%%%%%%%%%%%%%%%%%%%%%%%%%%%%%%%%%%%%%
\subsection{Spatial Sampling}
\label{subsec:spatial-sampling}
%%%%%%%%%%%%%%%%%%%%%%%%%%%%%%%%%%%%%%%%%%%%%%%%%%%%%%%%%%%%

Spatial dependence affects how locations should be sampled.

Possible designs include:

\[
\boxed{
\text{simple random spatial sampling},
}
\]

\[
\boxed{
\text{systematic grids},
}
\]

\[
\boxed{
\text{stratified spatial sampling},
}
\]

or

\[
\boxed{
\text{adaptive sampling}.
}
\]

%%%%%%%%%%%%%%%%%%%%%%%%%%%%%%%%%%%%%%%%%%%%%%%%%%%%%%%%%%%%
\subsection{Regular Grid Sampling}
\label{subsec:regular-grid-sampling}
%%%%%%%%%%%%%%%%%%%%%%%%%%%%%%%%%%%%%%%%%%%%%%%%%%%%%%%%%%%%

A regular grid provides uniform spatial coverage.

It is useful for mapping and interpolation.

However, systematic alignment with an unrecognized periodic spatial
pattern can create bias.

%%%%%%%%%%%%%%%%%%%%%%%%%%%%%%%%%%%%%%%%%%%%%%%%%%%%%%%%%%%%
\subsection{Spatially Balanced Sampling}
\label{subsec:spatially-balanced-sampling}
%%%%%%%%%%%%%%%%%%%%%%%%%%%%%%%%%%%%%%%%%%%%%%%%%%%%%%%%%%%%

Spatially balanced designs spread sample locations across the study
region.

The goal is to avoid excessive clustering of sampling sites.

This can improve coverage and prediction over a fixed sampling budget.

%%%%%%%%%%%%%%%%%%%%%%%%%%%%%%%%%%%%%%%%%%%%%%%%%%%%%%%%%%%%
\subsection{Preferential Sampling}
\label{subsec:preferential-sampling}
%%%%%%%%%%%%%%%%%%%%%%%%%%%%%%%%%%%%%%%%%%%%%%%%%%%%%%%%%%%%

Preferential sampling occurs when the locations selected for observation
depend on the underlying spatial process.

For example, monitors may be placed more frequently where pollution is
expected to be high.

Then the observed locations are not independent of the process being
measured.

Ignoring this mechanism can bias inference.

%%%%%%%%%%%%%%%%%%%%%%%%%%%%%%%%%%%%%%%%%%%%%%%%%%%%%%%%%%%%
\subsection{Spatial Interpolation}
\label{subsec:spatial-interpolation}
%%%%%%%%%%%%%%%%%%%%%%%%%%%%%%%%%%%%%%%%%%%%%%%%%%%%%%%%%%%%

Spatial interpolation predicts values at unsampled locations.

Methods include:

\[
\boxed{
\text{nearest-neighbor interpolation},
}
\]

\[
\boxed{
\text{inverse-distance weighting},
}
\]

\[
\boxed{
\text{splines},
}
\]

and

\[
\boxed{
\text{kriging}.
}
\]

Kriging differs by explicitly using a probabilistic spatial dependence
model.

%%%%%%%%%%%%%%%%%%%%%%%%%%%%%%%%%%%%%%%%%%%%%%%%%%%%%%%%%%%%
\subsection{Inverse-Distance Weighting}
\label{subsec:inverse-distance-weighting}
%%%%%%%%%%%%%%%%%%%%%%%%%%%%%%%%%%%%%%%%%%%%%%%%%%%%%%%%%%%%

Inverse-distance weighting predicts

\[
\boxed{
\widehat{Y}
(
\mathbf{s}_0
)
=
\frac{
\displaystyle
\sum_i
d_i^{-p}Y_i
}{
\displaystyle
\sum_i
d_i^{-p}
},
}
\]

where

\[
d_i
=
\|
\mathbf{s}_i-\mathbf{s}_0
\|.
\]

The exponent \(p>0\) controls how rapidly influence decreases with
distance.

This is deterministic rather than model-based spatial prediction.

%%%%%%%%%%%%%%%%%%%%%%%%%%%%%%%%%%%%%%%%%%%%%%%%%%%%%%%%%%%%
\subsection{Spatial Cross-Validation}
\label{subsec:spatial-cross-validation}
%%%%%%%%%%%%%%%%%%%%%%%%%%%%%%%%%%%%%%%%%%%%%%%%%%%%%%%%%%%%

Random cross-validation can be optimistic when nearby training and test
observations are strongly correlated.

Spatial cross-validation may instead withhold:

\[
\boxed{
\text{blocks},
}
\]

\[
\boxed{
\text{regions},
}
\]

or

\[
\boxed{
\text{spatially separated folds}.
}
\]

This better evaluates geographic generalization.

%%%%%%%%%%%%%%%%%%%%%%%%%%%%%%%%%%%%%%%%%%%%%%%%%%%%%%%%%%%%
\subsection{Modifiable Areal Unit Problem}
\label{subsec:maup}
%%%%%%%%%%%%%%%%%%%%%%%%%%%%%%%%%%%%%%%%%%%%%%%%%%%%%%%%%%%%

For areal data, conclusions can depend on how spatial regions are defined.

This is called the modifiable areal unit problem.

Changing:

\[
\boxed{
\text{boundary placement}
}
\]

or

\[
\boxed{
\text{aggregation scale}
}
\]

can change correlations, rates, and regression coefficients.

%%%%%%%%%%%%%%%%%%%%%%%%%%%%%%%%%%%%%%%%%%%%%%%%%%%%%%%%%%%%
\subsection{Ecological Fallacy}
\label{subsec:ecological-fallacy-spatial}
%%%%%%%%%%%%%%%%%%%%%%%%%%%%%%%%%%%%%%%%%%%%%%%%%%%%%%%%%%%%

Relationships observed among aggregated geographic regions need not hold
for individuals within those regions.

For example,

\[
\boxed{
\text{region-level association}
}
\]

does not automatically imply

\[
\boxed{
\text{individual-level association}.
}
\]

This is the ecological fallacy.

%%%%%%%%%%%%%%%%%%%%%%%%%%%%%%%%%%%%%%%%%%%%%%%%%%%%%%%%%%%%
\subsection{Change of Support}
\label{subsec:change-of-support}
%%%%%%%%%%%%%%%%%%%%%%%%%%%%%%%%%%%%%%%%%%%%%%%%%%%%%%%%%%%%

Spatial data may be observed at one scale but predictions are required at
another.

Examples include:

\[
\boxed{
\text{point measurements}
\rightarrow
\text{regional averages},
}
\]

or

\[
\boxed{
\text{county totals}
\rightarrow
\text{smaller areas}.
}
\]

This is called a change-of-support problem.

%%%%%%%%%%%%%%%%%%%%%%%%%%%%%%%%%%%%%%%%%%%%%%%%%%%%%%%%%%%%
\subsection{Spatiotemporal Data}
\label{subsec:spatiotemporal-data}
%%%%%%%%%%%%%%%%%%%%%%%%%%%%%%%%%%%%%%%%%%%%%%%%%%%%%%%%%%%%

Many datasets vary over both space and time.

Write

\[
\boxed{
Y(\mathbf{s},t).
}
\]

Examples include:

\[
\boxed{
\text{weather fields},
}
\]

\[
\boxed{
\text{air pollution},
}
\]

\[
\boxed{
\text{disease spread},
}
\]

or

\[
\boxed{
\text{ocean measurements}.
}
\]

%%%%%%%%%%%%%%%%%%%%%%%%%%%%%%%%%%%%%%%%%%%%%%%%%%%%%%%%%%%%
\subsection{Spatiotemporal Covariance}
\label{subsec:spatiotemporal-covariance}
%%%%%%%%%%%%%%%%%%%%%%%%%%%%%%%%%%%%%%%%%%%%%%%%%%%%%%%%%%%%

A spatiotemporal covariance function may be written

\[
\boxed{
C
(
\mathbf{h},
u
)
=
\operatorname{Cov}
[
Y(\mathbf{s},t),
Y(\mathbf{s}+\mathbf{h},t+u)
].
}
\]

Here,

\[
\mathbf{h}
\]

is spatial lag and

\[
u
\]

is temporal lag.

%%%%%%%%%%%%%%%%%%%%%%%%%%%%%%%%%%%%%%%%%%%%%%%%%%%%%%%%%%%%
\subsection{Separable Spatiotemporal Covariance}
\label{subsec:separable-spatiotemporal-covariance}
%%%%%%%%%%%%%%%%%%%%%%%%%%%%%%%%%%%%%%%%%%%%%%%%%%%%%%%%%%%%

A separable covariance has the form

\[
\boxed{
C
(
\mathbf{h},
u
)
=
C_S(\mathbf{h})
C_T(u).
}
\]

This simplifies modeling and computation.

However, real systems may have interactions between spatial and temporal
dependence that violate separability.

%%%%%%%%%%%%%%%%%%%%%%%%%%%%%%%%%%%%%%%%%%%%%%%%%%%%%%%%%%%%
\subsection{Spatial Prediction Uncertainty}
\label{subsec:spatial-prediction-uncertainty}
%%%%%%%%%%%%%%%%%%%%%%%%%%%%%%%%%%%%%%%%%%%%%%%%%%%%%%%%%%%%

Prediction uncertainty generally increases when:

\[
\boxed{
\text{prediction sites are far from observations},
}
\]

\[
\boxed{
\text{measurement noise is large},
}
\]

or

\[
\boxed{
\text{spatial dependence is weak}.
}
\]

Dense nearby observations usually reduce uncertainty, though redundancy
among highly correlated observations matters.

%%%%%%%%%%%%%%%%%%%%%%%%%%%%%%%%%%%%%%%%%%%%%%%%%%%%%%%%%%%%
\subsection{Spatial Model Diagnostics}
\label{subsec:spatial-model-diagnostics}
%%%%%%%%%%%%%%%%%%%%%%%%%%%%%%%%%%%%%%%%%%%%%%%%%%%%%%%%%%%%

After fitting a spatial model, one should examine:

\[
\boxed{
\text{residual maps},
}
\]

\[
\boxed{
\text{residual spatial autocorrelation},
}
\]

\[
\boxed{
\text{variogram behavior},
}
\]

\[
\boxed{
\text{prediction errors},
}
\]

and

\[
\boxed{
\text{spatial cross-validation}.
}
\]

Residual spatial structure suggests the model has not captured all
dependence.

%%%%%%%%%%%%%%%%%%%%%%%%%%%%%%%%%%%%%%%%%%%%%%%%%%%%%%%%%%%%
\subsection{Spatial Statistical Workflow}
\label{subsec:spatial-statistical-workflow}
%%%%%%%%%%%%%%%%%%%%%%%%%%%%%%%%%%%%%%%%%%%%%%%%%%%%%%%%%%%%

A practical spatial workflow is:

\begin{enumerate}
    \item identify the spatial data type;
    \item map the observations;
    \item verify coordinate systems;
    \item inspect spatial coverage;
    \item identify relevant covariates;
    \item examine first-order spatial trends;
    \item examine pairwise distances or neighborhood structure;
    \item inspect empirical spatial dependence;
    \item select a covariance, variogram, or weights structure;
    \item fit the spatial model;
    \item examine residual spatial autocorrelation;
    \item quantify spatial prediction uncertainty;
    \item validate using spatially appropriate holdout procedures;
    \item assess sensitivity to neighborhood or covariance choices;
    \item interpret results at the correct spatial scale.
\end{enumerate}

%%%%%%%%%%%%%%%%%%%%%%%%%%%%%%%%%%%%%%%%%%%%%%%%%%%%%%%%%%%%
\subsection{Choosing a Spatial Method}
\label{subsec:choosing-spatial-method}
%%%%%%%%%%%%%%%%%%%%%%%%%%%%%%%%%%%%%%%%%%%%%%%%%%%%%%%%%%%%

For continuous measurements over a region:

\[
\boxed{
\text{geostatistical models and kriging}.
}
\]

For regional counts or rates:

\[
\boxed{
\text{areal spatial models}.
}
\]

For event locations:

\[
\boxed{
\text{point-process methods}.
}
\]

For residual spatial dependence in regression:

\[
\boxed{
\text{spatial regression}.
}
\]

For spatially correlated latent effects:

\[
\boxed{
\text{Gaussian-process, CAR, or SAR models}.
}
\]

%%%%%%%%%%%%%%%%%%%%%%%%%%%%%%%%%%%%%%%%%%%%%%%%%%%%%%%%%%%%
\subsection{Common Errors}
\label{subsec:spatial-common-errors}
%%%%%%%%%%%%%%%%%%%%%%%%%%%%%%%%%%%%%%%%%%%%%%%%%%%%%%%%%%%%

\subsubsection{Treating Spatial Observations as Independent}

Nearby observations may contain redundant information.

Ignoring spatial dependence can lead to underestimated standard errors
and overconfident inference.

%%%%%%%%%%%%%%%%%%%%%%%%%%%%%%%%%%%%%%%%%%%%%%%%%%%%%%%%%%%%
\subsubsection{Using Latitude and Longitude as Ordinary Euclidean Coordinates}

Large geographic regions require appropriate geographic distance or map
projection methods.

%%%%%%%%%%%%%%%%%%%%%%%%%%%%%%%%%%%%%%%%%%%%%%%%%%%%%%%%%%%%
\subsubsection{Confusing Spatial Trend with Spatial Correlation}

A broad spatial gradient can create apparent correlation.

First-order trend should be separated from second-order dependence when
possible.

%%%%%%%%%%%%%%%%%%%%%%%%%%%%%%%%%%%%%%%%%%%%%%%%%%%%%%%%%%%%
\subsubsection{Assuming Isotropy Without Checking Direction}

Dependence may differ by direction because of terrain, wind, water flow,
roads, or other mechanisms.

%%%%%%%%%%%%%%%%%%%%%%%%%%%%%%%%%%%%%%%%%%%%%%%%%%%%%%%%%%%%
\subsubsection{Interpreting the Empirical Variogram Without Enough Pairs}

Semivariogram estimates at sparsely populated distance bins can be highly
variable.

The number of pairs should be considered.

%%%%%%%%%%%%%%%%%%%%%%%%%%%%%%%%%%%%%%%%%%%%%%%%%%%%%%%%%%%%
\subsubsection{Treating the Nugget as Only Measurement Error}

A nugget can also represent unresolved variation occurring at scales
smaller than the sampling distance.

%%%%%%%%%%%%%%%%%%%%%%%%%%%%%%%%%%%%%%%%%%%%%%%%%%%%%%%%%%%%
\subsubsection{Assuming Kriging Is Just Distance Weighting}

Kriging accounts for covariance among observations, not only distance
from the target location.

%%%%%%%%%%%%%%%%%%%%%%%%%%%%%%%%%%%%%%%%%%%%%%%%%%%%%%%%%%%%
\subsubsection{Ignoring Covariance-Model Uncertainty}

Kriging variance often treats estimated covariance parameters as known.

Actual uncertainty may therefore be larger.

%%%%%%%%%%%%%%%%%%%%%%%%%%%%%%%%%%%%%%%%%%%%%%%%%%%%%%%%%%%%
\subsubsection{Choosing a Spatial Weights Matrix Arbitrarily}

Moran statistics and spatial regression results can depend strongly on
the definition of \(W\).

Neighborhood structure should have scientific justification.

%%%%%%%%%%%%%%%%%%%%%%%%%%%%%%%%%%%%%%%%%%%%%%%%%%%%%%%%%%%%
\subsubsection{Interpreting Moran's \(I\) as a Correlation Coefficient}

Moran's \(I\) is related to spatial autocorrelation but is not bounded in
exactly the same way as an ordinary Pearson correlation in all settings.

%%%%%%%%%%%%%%%%%%%%%%%%%%%%%%%%%%%%%%%%%%%%%%%%%%%%%%%%%%%%
\subsubsection{Ignoring Multiple Testing in Local Spatial Statistics}

A map containing many local significance tests can produce false
clusters by chance.

Multiplicity or simulation-based reference procedures should be
considered.

%%%%%%%%%%%%%%%%%%%%%%%%%%%%%%%%%%%%%%%%%%%%%%%%%%%%%%%%%%%%
\subsubsection{Confusing Spatial Lag and Spatial Error Models}

They represent different forms of dependence.

One cannot generally substitute one for the other without changing the
scientific interpretation.

%%%%%%%%%%%%%%%%%%%%%%%%%%%%%%%%%%%%%%%%%%%%%%%%%%%%%%%%%%%%
\subsubsection{Ignoring Spatial Confounding}

A spatially smooth predictor may overlap strongly with a latent spatial
effect.

This can make regression coefficients unstable or difficult to
interpret.

%%%%%%%%%%%%%%%%%%%%%%%%%%%%%%%%%%%%%%%%%%%%%%%%%%%%%%%%%%%%
\subsubsection{Using Random Cross-Validation for Strongly Spatial Data}

Nearby training observations may make predictions appear unrealistically
accurate.

Spatial blocking gives a more demanding generalization assessment.

%%%%%%%%%%%%%%%%%%%%%%%%%%%%%%%%%%%%%%%%%%%%%%%%%%%%%%%%%%%%
\subsubsection{Ignoring Edge Effects in Point Patterns}

Events near study boundaries have fewer observable neighboring locations.

Uncorrected neighborhood summaries may therefore be biased.

%%%%%%%%%%%%%%%%%%%%%%%%%%%%%%%%%%%%%%%%%%%%%%%%%%%%%%%%%%%%
\subsubsection{Equating Clustering with a Point-Interaction Process}

Apparent clustering may arise simply because intensity varies across
space.

First-order inhomogeneity must be separated from interaction.

%%%%%%%%%%%%%%%%%%%%%%%%%%%%%%%%%%%%%%%%%%%%%%%%%%%%%%%%%%%%
\subsubsection{Inferring Individual Relationships from Areal Data}

Area-level associations do not automatically hold at the individual
level.

%%%%%%%%%%%%%%%%%%%%%%%%%%%%%%%%%%%%%%%%%%%%%%%%%%%%%%%%%%%%
\subsection{Applications}
\label{subsec:spatial-applications}
%%%%%%%%%%%%%%%%%%%%%%%%%%%%%%%%%%%%%%%%%%%%%%%%%%%%%%%%%%%%

\begin{applicationbox}

\textbf{Environmental Science.}
Spatial statistics models dissolved oxygen, temperature, salinity,
pollution, soil properties, and ecological measurements across
monitoring locations.

\medskip

\textbf{Epidemiology.}
Disease incidence, environmental exposures, and regional health outcomes
may exhibit geographic clustering and spatial dependence.

\medskip

\textbf{Ecology.}
Species abundance, habitat suitability, vegetation, and point patterns of
organisms are naturally spatial.

\medskip

\textbf{Geology.}
Mineral concentrations, subsurface properties, and geological fields are
frequently interpolated using geostatistical methods.

\medskip

\textbf{Meteorology.}
Temperature, rainfall, wind, and pressure are spatial fields that also
change through time.

\medskip

\textbf{Agriculture.}
Crop yield, soil nutrients, irrigation response, and pest occurrence can
be modeled over fields and regions.

\medskip

\textbf{Economics and Social Science.}
Income, employment, housing values, and demographic variables may exhibit
dependence across neighboring geographic regions.

\medskip

\textbf{Engineering.}
Spatial sensor networks, structural monitoring, and environmental
measurement systems require prediction over unsampled locations.

\end{applicationbox}

%%%%%%%%%%%%%%%%%%%%%%%%%%%%%%%%%%%%%%%%%%%%%%%%%%%%%%%%%%%%
\subsection{Exercises}
\label{subsec:spatial-exercises}
%%%%%%%%%%%%%%%%%%%%%%%%%%%%%%%%%%%%%%%%%%%%%%%%%%%%%%%%%%%%

\begin{exercise}[1]
Define geostatistical data.
\end{exercise}

\begin{exercise}[1]
Define areal data.
\end{exercise}

\begin{exercise}[1]
Define a spatial point pattern.
\end{exercise}

\begin{exercise}[1]
Define spatial autocorrelation.
\end{exercise}

\begin{exercise}[1]
Define isotropy.
\end{exercise}

\begin{exercise}[1]
Define anisotropy.
\end{exercise}

\begin{exercise}[1]
Define a semivariogram.
\end{exercise}

\begin{exercise}[1]
Define nugget, sill, and range.
\end{exercise}

\begin{exercise}[1]
Define kriging.
\end{exercise}

\begin{exercise}[1]
Define a spatial weights matrix.
\end{exercise}

\begin{exercise}[1]
Define Moran's \(I\).
\end{exercise}

\begin{exercise}[1]
Define an intensity function for a point process.
\end{exercise}

\begin{exercise}[2]
Compute the Euclidean distance between

\[
(1,2)
\]

and

\[
(4,6).
\]
\end{exercise}

\begin{exercise}[2]
Suppose

\[
C(0)=10
\]

and

\[
C(h)=4.
\]

Compute the semivariogram value \(\gamma(h)\).
\end{exercise}

\begin{exercise}[2]
Suppose four location pairs in a distance bin have differences

\[
1,\ -2,\ 3,\ -4.
\]

Compute the empirical semivariance.
\end{exercise}

\begin{exercise}[2]
For

\[
C(h)
=
5e^{-h/2},
\]

compute \(C(0)\).
\end{exercise}

\begin{exercise}[2]
For the same covariance model, compute \(C(2)\).
\end{exercise}

\begin{exercise}[2]
Suppose a row of a spatial weights matrix is

\[
(0,1,1,2).
\]

Row-standardize the weights.
\end{exercise}

\begin{exercise}[2]
Suppose a homogeneous Poisson point process has intensity

\[
\lambda=3
\]

events per square kilometer.

Find the expected number of events in a region of area \(5\).
\end{exercise}

\begin{exercise}[2]
Under the same process, state the distribution of the event count in the
region.
\end{exercise}

\begin{exercise}[2]
Under complete spatial randomness, compute

\[
K(2).
\]
\end{exercise}

\begin{exercise}[2]
If an estimated \(K(2)\) greatly exceeds the CSR value, what type of
pattern may be present?
\end{exercise}

\begin{exercise}[3]
Show that

\[
\gamma(\mathbf{h})
=
C(0)-C(\mathbf{h})
\]

under second-order stationarity.
\end{exercise}

\begin{exercise}[3]
Explain why a valid covariance function must produce a positive
semidefinite covariance matrix.
\end{exercise}

\begin{exercise}[3]
Explain why isotropy is stronger than second-order stationarity.
\end{exercise}

\begin{exercise}[3]
Compare exponential and Gaussian covariance models conceptually.
\end{exercise}

\begin{exercise}[3]
Explain how directional variograms can reveal anisotropy.
\end{exercise}

\begin{exercise}[3]
Derive the simple kriging weights from minimum mean squared prediction
error.
\end{exercise}

\begin{exercise}[3]
Derive the simple kriging prediction variance.
\end{exercise}

\begin{exercise}[3]
Explain why ordinary kriging imposes

\[
\sum_i\lambda_i=1.
\]
\end{exercise}

\begin{exercise}[3]
Explain the difference between ordinary and universal kriging.
\end{exercise}

\begin{exercise}[3]
Show how kriging follows from the conditional distribution of a
multivariate normal vector.
\end{exercise}

\begin{exercise}[3]
Explain why spatially close observations may be statistically redundant.
\end{exercise}

\begin{exercise}[3]
Explain how row-standardization changes interpretation of a spatial lag.
\end{exercise}

\begin{exercise}[3]
Explain the difference between Moran's \(I\) and Geary's \(C\).
\end{exercise}

\begin{exercise}[3]
Explain the difference between global Moran's \(I\) and local Moran
statistics.
\end{exercise}

\begin{exercise}[3]
Explain why multiple-testing correction may be needed for LISA maps.
\end{exercise}

\begin{exercise}[3]
Explain the difference between spatial lag and spatial error regression.
\end{exercise}

\begin{exercise}[3]
Derive the reduced form

\[
\mathbf{y}
=
(
I-\rho W
)^{-1}
X\boldsymbol{\beta}
+
(
I-\rho W
)^{-1}
\boldsymbol{\varepsilon}.
\]
\end{exercise}

\begin{exercise}[3]
Explain how a CAR prior encourages neighboring effects to be similar.
\end{exercise}

\begin{exercise}[3]
Explain the conceptual difference between CAR and SAR models.
\end{exercise}

\begin{exercise}[3]
Explain complete spatial randomness.
\end{exercise}

\begin{exercise}[3]
Explain how inhomogeneous intensity can create apparent point clustering.
\end{exercise}

\begin{exercise}[3]
Explain why edge correction is needed for nearest-neighbor or
\(K\)-function analysis.
\end{exercise}

\begin{exercise}[3]
Compare inverse-distance weighting and kriging.
\end{exercise}

\begin{exercise}[3]
Explain why spatial cross-validation can be more realistic than random
cross-validation.
\end{exercise}

\begin{exercise}[3]
Explain the modifiable areal unit problem.
\end{exercise}

\begin{exercise}[3]
Explain the ecological fallacy.
\end{exercise}

\begin{exercise}[3]
Explain the concept of preferential sampling.
\end{exercise}

\begin{exercise}[3]
Explain the meaning of a separable spatiotemporal covariance function.
\end{exercise}

\begin{exercise}[4]
Study conditions for validity of isotropic covariance functions.
\end{exercise}

\begin{exercise}[4]
Derive properties of the Matérn covariance family.
\end{exercise}

\begin{exercise}[4]
Study intrinsic random functions and generalized covariance models.
\end{exercise}

\begin{exercise}[4]
Derive ordinary kriging using Lagrange multipliers.
\end{exercise}

\begin{exercise}[4]
Derive universal kriging under a linear mean model.
\end{exercise}

\begin{exercise}[4]
Study restricted maximum likelihood estimation of spatial covariance
parameters.
\end{exercise}

\begin{exercise}[4]
Investigate covariance tapering for large spatial datasets.
\end{exercise}

\begin{exercise}[4]
Study Gaussian-process approximation methods.
\end{exercise}

\begin{exercise}[4]
Derive the expectation and variance of Moran's \(I\) under a randomization
null.
\end{exercise}

\begin{exercise}[4]
Study permutation tests for global and local spatial autocorrelation.
\end{exercise}

\begin{exercise}[4]
Derive the covariance structure of a SAR model.
\end{exercise}

\begin{exercise}[4]
Study propriety and identifiability of intrinsic CAR models.
\end{exercise}

\begin{exercise}[4]
Investigate Bayesian disease-mapping models.
\end{exercise}

\begin{exercise}[4]
Study spatial generalized linear mixed models.
\end{exercise}

\begin{exercise}[4]
Derive likelihoods for homogeneous Poisson point processes.
\end{exercise}

\begin{exercise}[4]
Study inhomogeneous Poisson point processes.
\end{exercise}

\begin{exercise}[4]
Investigate Cox and log-Gaussian Cox processes.
\end{exercise}

\begin{exercise}[4]
Derive Ripley's \(K\)-function under complete spatial randomness.
\end{exercise}

\begin{exercise}[4]
Study edge corrections for point-process summary statistics.
\end{exercise}

\begin{exercise}[4]
Investigate spatial point-process interaction models.
\end{exercise}

\begin{exercise}[4]
Study preferential sampling models.
\end{exercise}

\begin{exercise}[4]
Investigate nonseparable spatiotemporal covariance functions.
\end{exercise}

\begin{exercise}[4]
Study spatiotemporal Gaussian processes.
\end{exercise}

%%%%%%%%%%%%%%%%%%%%%%%%%%%%%%%%%%%%%%%%%%%%%%%%%%%%%%%%%%%%
\subsection{Conceptual Summary}
\label{subsec:spatial-summary}
%%%%%%%%%%%%%%%%%%%%%%%%%%%%%%%%%%%%%%%%%%%%%%%%%%%%%%%%%%%%

Spatial statistics studies observations indexed by location:

\[
\boxed{
Y(\mathbf{s}).
}
\]

Second-order spatial dependence is described by

\[
\boxed{
C(\mathbf{h})
=
\operatorname{Cov}
[
Y(\mathbf{s}),
Y(\mathbf{s}+\mathbf{h})
].
}
\]

For isotropic processes,

\[
\boxed{
C(\mathbf{h})
=
C
(
\|\mathbf{h}\|
).
}
\]

The semivariogram is

\[
\boxed{
\gamma(\mathbf{h})
=
\frac12
\operatorname{Var}
[
Y(\mathbf{s}+\mathbf{h})
-
Y(\mathbf{s})
].
}
\]

Under stationarity,

\[
\boxed{
\gamma(\mathbf{h})
=
C(0)-C(\mathbf{h}).
}
\]

Kriging predicts

\[
\boxed{
\widehat{Y}
(
\mathbf{s}_0
)
=
\sum_i
\lambda_i
Y
(
\mathbf{s}_i
)
}
\]

using weights based on spatial dependence.

For areal data, dependence may be represented through a spatial weights
matrix

\[
\boxed{
W.
}
\]

A global measure of spatial autocorrelation is Moran's statistic

\[
\boxed{
I
=
\frac{
n
}{
S_0
}
\frac{
\sum_i\sum_j
w_{ij}
(
y_i-\overline{y}
)
(
y_j-\overline{y}
)
}{
\sum_i
(
y_i-\overline{y}
)^2
}.
}
\]

For spatial point patterns, a homogeneous Poisson process satisfies

\[
\boxed{
N(A)
\sim
\operatorname{Poisson}
(
\lambda|A|
).
}
\]

Under complete spatial randomness in two dimensions,

\[
\boxed{
K(r)
=
\pi r^2.
}
\]

Spatial statistics therefore combines

\[
\boxed{
\text{geometry}
+
\text{probability}
+
\text{covariance}
+
\text{regression}
+
\text{stochastic processes}.
}
\]

%%%%%%%%%%%%%%%%%%%%%%%%%%%%%%%%%%%%%%%%%%%%%%%%%%%%%%%%%%%%
\subsection{Section Connections}
\label{subsec:spatial-connections}
%%%%%%%%%%%%%%%%%%%%%%%%%%%%%%%%%%%%%%%%%%%%%%%%%%%%%%%%%%%%

\chapterconnection{
Spatial Statistics extends the dependence ideas of
Section~\ref{sec:time-series-analysis} from temporal indexing to spatial
locations. Geostatistical models use covariance and variogram structures
to describe continuously indexed spatial fields, while areal models use
neighborhood structures and spatial weights matrices. Point-process
models instead treat event locations themselves as random. Gaussian
random fields connect spatial statistics with multivariate normal theory,
while spatial regression extends the models of
Section~\ref{sec:regression} to geographically dependent observations.
Spatiotemporal models combine both spatial and temporal dependence. The
next section, Survival Analysis, studies time-to-event outcomes,
censoring, hazard functions, and statistical models for event times.
}

%%%%%%%%%%%%%%%%%%%%%%%%%%%%%%%%%%%%%%%%%%%%%%%%%%%%%%%%%%%%
% SECTION SOURCE TRACKING
%%%%%%%%%%%%%%%%%%%%%%%%%%%%%%%%%%%%%%%%%%%%%%%%%%%%%%%%%%%%

%------------------------------------------------------------
% 8.16 Spatial Statistics
%------------------------------------------------------------
%
% Source basis:
%   - Standard spatial statistical theory.
%   - Standard geostatistical methods.
%   - Standard areal and point-process models.
%
% Used for:
%   - spatial data types;
%   - geographic coordinates and distance;
%   - first- and second-order spatial structure;
%   - spatial stationarity;
%   - isotropy and anisotropy;
%   - covariance and correlation functions;
%   - exponential, Gaussian, and Matern covariance models;
%   - intrinsic stationarity;
%   - semivariograms and variograms;
%   - empirical semivariograms;
%   - nugget, sill, and range;
%   - spherical and exponential variogram models;
%   - directional variograms;
%   - kriging;
%   - simple, ordinary, and universal kriging;
%   - kriging variance;
%   - Gaussian random fields;
%   - spatial weights matrices;
%   - adjacency and distance weights;
%   - spatial lags;
%   - Moran's I;
%   - Geary's C;
%   - local spatial association;
%   - spatial regression;
%   - spatial lag models;
%   - spatial error models;
%   - spatial confounding;
%   - CAR and SAR models;
%   - spatial generalized linear models;
%   - point processes;
%   - intensity functions;
%   - homogeneous Poisson point processes;
%   - complete spatial randomness;
%   - clustering and inhibition;
%   - Ripley's K-function;
%   - edge effects;
%   - inhomogeneous point processes;
%   - Cox-process previews;
%   - spatial sampling;
%   - preferential sampling;
%   - spatial interpolation;
%   - spatial cross-validation;
%   - modifiable areal unit problems;
%   - ecological fallacy;
%   - change of support;
%   - spatiotemporal covariance;
%   - model diagnostics;
%   - applications and exercises.
%
% Notes:
%   - Survival Analysis is developed next in Section 8.17.
%   - Time Series Analysis provides the parallel theory for temporal
%     dependence.
%   - Multivariate Statistics provides covariance and Gaussian theory.
%   - Statistical Learning later develops Gaussian processes and
%     spatially structured predictive models further.
%
%%%%%%%%%%%%%%%%%%%%%%%%%%%%%%%%%%%%%%%%%%%%%%%%%%%%%%%%%%%%